\chapter{Communication Complexity --- Bit String}
%%% generated by the blueprint pipeline from TCSlib.CommunicationComplexity.DeterministicCC.BitString
\section{Overview}

This module introduces $n$-bit strings as Boolean-valued functions on $\mathrm{Fin}\,n$
and develops the \emph{signed inner product} of two such strings under the standard
$\{0,1\} \to \{\pm 1\}$ correspondence.  The main results relate the signed inner
product to the Hamming distance: it equals $n - 2\,\mathrm{hammingDist}(x,y)$, and it
is additive under concatenation and multiplicative under repetition.

%%%%%%%%%%%%%%%%%%%%%%%%%%%%%%%%%%%%%%%%%%%%%%%%%%%%%%%%%%%%%%%%%%%%%%%%%%%%%%%
\section{Declarations}

\begin{definition}[Bit string]
\lean{CommunicationComplexity.BitString}
\label{CommunicationComplexity.BitString}
\leanok
An $n$-bit string is a function $x : \mathrm{Fin}\,n \to \mathrm{Bool}$.
The type $\texttt{CommunicationComplexity.BitString}\,n$ is an abbreviation for
$\mathrm{Fin}\,n \to \mathrm{Bool}$.
\end{definition}

\begin{definition}[Signed inner product]
\lean{CommunicationComplexity.BitString.signedInner}
\label{CommunicationComplexity.BitString.signedInner}
\leanok
\uses{CommunicationComplexity.BitString}
Given two $n$-bit strings $x, y : \mathrm{Fin}\,n \to \mathrm{Bool}$, their
\emph{signed inner product} is the integer
\[
  \mathrm{signedInner}(x, y) \;=\; \sum_{i=0}^{n-1}
    \begin{cases} 1 & \text{if } x_i = y_i,\\ -1 & \text{if } x_i \neq y_i.\end{cases}
\]
Each coordinate where $x$ and $y$ agree contributes $+1$, and each disagreeing
coordinate contributes $-1$.
\end{definition}

\begin{definition}[Agreement count]
\lean{CommunicationComplexity.BitString.agreementCount}
\label{CommunicationComplexity.BitString.agreementCount}
\leanok
\uses{CommunicationComplexity.BitString}
The \emph{agreement count} of two $n$-bit strings $x$ and $y$ is the number of
coordinates $i \in \mathrm{Fin}\,n$ on which $x_i = y_i$, i.e.,
$\#\{i : x_i = y_i\}$ as a natural number.
\end{definition}

\begin{theorem}[Agreement count plus Hamming distance equals length]
\lean{CommunicationComplexity.BitString.agreementCount_add_hammingDist_eq_length}
\label{CommunicationComplexity.BitString.agreementCount_add_hammingDist_eq_length}
\leanok
\uses{CommunicationComplexity.BitString, CommunicationComplexity.BitString.agreementCount}
\difficulty{2}
Let $x$ and $y$ be $n$-bit strings. Then the number of coordinates on which $x$ and $y$
agree, together with their Hamming distance $d(x, y)$ — the number of coordinates on
which they differ — sums to the length:
\[
  \#\{i : x_i = y_i\} + d(x, y) = n.
\]
\end{theorem}

\begin{theorem}[Agreement count as complement of Hamming distance]
\lean{CommunicationComplexity.BitString.agreementCount_eq_length_sub_hammingDist}
\label{CommunicationComplexity.BitString.agreementCount_eq_length_sub_hammingDist}
\leanok
\uses{CommunicationComplexity.BitString, CommunicationComplexity.BitString.agreementCount, CommunicationComplexity.BitString.agreementCount_add_hammingDist_eq_length}
\difficulty{2}
Let $x, y : \mathrm{Fin}\,n \to \mathrm{Bool}$ be two $n$-bit strings, with agreement
count $\mathrm{agr}(x,y)$ the number of coordinates on which they coincide and
$d_H(x,y)$ their Hamming distance (the number of coordinates on which they differ).
Then, as an identity of integers,
\[
  \mathrm{agr}(x,y) = n - d_H(x,y).
\]
\end{theorem}

\begin{theorem}[Signed inner product via Hamming distance]
\lean{CommunicationComplexity.BitString.signedInner_eq_length_sub_twice_hammingDist}
\label{CommunicationComplexity.BitString.signedInner_eq_length_sub_twice_hammingDist}
\leanok
\uses{CommunicationComplexity.BitString, CommunicationComplexity.BitString.signedInner, CommunicationComplexity.BitString.signedInner_append}
\difficulty{4}
Let $x, y$ be two $n$-bit strings, that is, functions from $\{0, 1, \dots, n-1\}$ to
$\{0,1\}$, and let $d(x,y)$ denote their Hamming distance, the number of coordinates on
which they disagree. Then their signed inner product — the sum over all coordinates of
$+1$ where $x$ and $y$ agree and $-1$ where they disagree — satisfies
\[
  \langle x, y\rangle_{\pm} \;=\; n - 2\,d(x, y).
\]
\end{theorem}

\begin{theorem}[Signed inner product as agreements minus Hamming distance]
\lean{CommunicationComplexity.BitString.signedInner_eq_agreementCount_sub_hammingDist}
\label{CommunicationComplexity.BitString.signedInner_eq_agreementCount_sub_hammingDist}
\leanok
\uses{CommunicationComplexity.BitString, CommunicationComplexity.BitString.agreementCount, CommunicationComplexity.BitString.agreementCount_eq_length_sub_hammingDist, CommunicationComplexity.BitString.signedInner, CommunicationComplexity.BitString.signedInner_eq_length_sub_twice_hammingDist}
\difficulty{2}
Let $x$ and $y$ be two $n$-bit strings. Their signed inner product $\langle x, y\rangle$
is the integer $\sum_{i} \varepsilon_i$, where $\varepsilon_i = +1$ at each coordinate
$i$ on which $x_i = y_i$ and $\varepsilon_i = -1$ at each coordinate on which $x_i \neq
y_i$. Writing $a(x,y)$ for the agreement count of $x$ and $y$, the number of coordinates
on which they agree, and $d_H(x,y)$ for the Hamming distance between them, the number of
coordinates on which they differ, one has
\[
  \langle x, y\rangle \;=\; a(x,y) - d_H(x,y),
\]
an equality of integers in which the two counts are regarded as integers.
\end{theorem}

\begin{theorem}[Additivity of the signed inner product under concatenation]
\lean{CommunicationComplexity.BitString.signedInner_append}
\label{CommunicationComplexity.BitString.signedInner_append}
\leanok
\uses{CommunicationComplexity.BitString.signedInner}
\difficulty{3}
Let $x_1, y_1 : \mathrm{Fin}\,m \to \mathrm{Bool}$ be $m$-bit strings and $x_2, y_2 :
\mathrm{Fin}\,n \to \mathrm{Bool}$ be $n$-bit strings, and write $x_1 \frown x_2$ for
the length-$(m+n)$ string obtained by concatenation, whose first $m$ coordinates are
those of $x_1$ and whose last $n$ coordinates are those of $x_2$. Writing $\langle x,
y\rangle$ for the signed inner product of two equal-length bit strings — the number of
coordinates on which they agree minus the number on which they disagree — the signed
inner product of the two concatenations splits as the sum of the block-wise signed inner
products:
\[
  \langle x_1 \frown x_2,\; y_1 \frown y_2\rangle
  \;=\; \langle x_1, y_1\rangle + \langle x_2, y_2\rangle.
\]
\end{theorem}

\begin{theorem}[Invariance of the signed inner product under reindexing]
\lean{CommunicationComplexity.BitString.signedInner_comp_cast}
\label{CommunicationComplexity.BitString.signedInner_comp_cast}
\leanok
\uses{CommunicationComplexity.BitString.signedInner}
\difficulty{1}
Let $m$ and $n$ be natural numbers with $m = n$, and let $x, y \colon \mathrm{Fin}\,n
\to \mathrm{Bool}$ be two $n$-bit strings, writing $\langle x, y\rangle$ for their
signed inner product. The equality $m = n$ gives a bijection $\sigma \colon
\mathrm{Fin}\,m \to \mathrm{Fin}\,n$, and precomposing with it turns $x$ and $y$ into
the $m$-bit strings $x \circ \sigma$ and $y \circ \sigma$. Then reindexing both strings
in this way leaves the signed inner product unchanged:
\[
  \langle x \circ \sigma,\; y \circ \sigma\rangle \;=\; \langle x, y\rangle.
\]
\end{theorem}

\begin{theorem}[Signed inner product scales under repetition]
\lean{CommunicationComplexity.BitString.signedInner_amplify}
\label{CommunicationComplexity.BitString.signedInner_amplify}
\leanok
\uses{CommunicationComplexity.BitString, CommunicationComplexity.BitString.signedInner, CommunicationComplexity.BitString.signedInner_append, CommunicationComplexity.BitString.signedInner_comp_cast}
\difficulty{4}
Let $x, y \colon \{0, \dots, n-1\} \to \{0,1\}$ be two $n$-bit strings, and write
$\langle x, y\rangle$ for their signed inner product, the integer $\sum_{i} \sigma_i$
where $\sigma_i = +1$ when $x_i = y_i$ and $\sigma_i = -1$ otherwise. For a natural
number $a$, let $x^{(a)}$ and $y^{(a)}$ denote the $an$-bit strings obtained by
concatenating $a$ consecutive copies of $x$ and of $y$, respectively. Then
\[
  \langle x^{(a)}, y^{(a)}\rangle \;=\; a \cdot \langle x, y\rangle.
\]
\end{theorem}
